% GENERATED FILE. Edit canonical lessons, then run npm run generate:books.
% canonical-source-hash: fnv1a64:3db19dd5476a32eb
% canonical-lessons: GU-C50-happiness, GU-C50-anger, GU-C50-fear, GU-C50-worry, GU-C50-tiredness

\chapter{How It Feels}
\label{ch:gu-how-it-feels}
\hlchaptermodality{50}

\begin{chapteropening}
\textbf{By the end of this chapter:} \emph{I can name happiness, anger, fear, worry and tiredness.}

Complete the last lesson of chapter 50: I can name happiness, anger, fear, worry and tiredness.
\end{chapteropening}

\section[\texorpdfstring{\gu{ખુશી} (khuśī)}{ખુશી (khuśī)}]{\gu{ખુશી} (khuśī) — happiness, joy}

\label{lesson:GU-C50-happiness}



\begin{quote}
Before the new one: say the Gujarati for hunger, then the Gujarati for thirst.
\end{quote}

\subsection*{You'll want to know: \gu{ખુશી}}
\textbf{\gu{ખુશી}} (\emph{khuśī}) — \textquotedblleft{}happiness, joy\textquotedblright{}.

\begin{cousinweb}
\textbf{\gu{ખુશી}} (\emph{khuśī}) is a loan, usually traced to Persian \emph{khushī}, \textquotedblleft{}happiness\textquotedblright{}, built on \emph{khush}, \textquotedblleft{}happy, pleasant\textquotedblright{}.
\end{cousinweb}

\begin{grammarlens}[title={What you've built}]
One more word for how it feels.
\end{grammarlens}

\subsection*{Your turn}
\begin{itemize}
\raggedright
  \item \emph{Say it:} \emph{khuśī}
  \item \emph{Say it:} \emph{khuśī}, once more
  \item \emph{Say it:} say \emph{khuśī}
  \item \emph{Recall:} say the Gujarati for hunger, then the Gujarati for thirst, then say \emph{khuśī} again
\end{itemize}

\begin{tcolorbox}[breakable,colback=teal!4,colframe=teal!35!black,title={Before you move on}]
What does \gu{ખુશી} mean? (\textquotedblleft{}Happiness, joy\textquotedblright{}.) Say it once more.
\end{tcolorbox}
\section[\texorpdfstring{\gu{ગુસ્સો} (gusso)}{ગુસ્સો (gusso)}]{\gu{ગુસ્સો} (gusso) — anger}

\label{lesson:GU-C50-anger}



\begin{quote}
Before the new one: say the Gujarati for thirst, then the Gujarati for happiness.
\end{quote}

\subsection*{You'll want to know: \gu{ગુસ્સો}}
\textbf{\gu{ગુસ્સો}} (\emph{gusso}) — \textquotedblleft{}anger\textquotedblright{}.

\begin{cousinweb}
\textbf{\gu{ગુસ્સો}} (\emph{gusso}) is a loan from Arabic \emph{ghuṣṣa}, \textquotedblleft{}a choking; grief\textquotedblright{}, which reached Gujarati through Persian. Its Arabic root is \emph{gh-ṣ-ṣ} (Lane).
\end{cousinweb}

\begin{grammarlens}[title={What you've built}]
One more word for how it feels.
\end{grammarlens}

\subsection*{Your turn}
\begin{itemize}
\raggedright
  \item \emph{Say it:} \emph{gusso}
  \item \emph{Say it:} \emph{gusso}, once more
  \item \emph{Say it:} say \emph{gusso}
  \item \emph{Recall:} say the Gujarati for thirst, then the Gujarati for happiness, then say \emph{gusso} again
\end{itemize}

\begin{tcolorbox}[breakable,colback=teal!4,colframe=teal!35!black,title={Before you move on}]
What does \gu{ગુસ્સો} mean? (\textquotedblleft{}Anger\textquotedblright{}.) Say it once more.
\end{tcolorbox}
\section[\texorpdfstring{\gu{બીક} (bīk)}{બીક (bīk)}]{\gu{બીક} (bīk) — fear}

\label{lesson:GU-C50-fear}



\begin{quote}
Before the new one: say the Gujarati for happiness, then the Gujarati for anger.
\end{quote}

\subsection*{You'll want to know: \gu{બીક}}
\textbf{\gu{બીક}} (\emph{bīk}) — \textquotedblleft{}fear\textquotedblright{}.

\begin{cousinweb}
\textbf{\gu{બીક}} (\emph{bīk}) belongs with the verb \emph{bīvũ}, \textquotedblleft{}to fear\textquotedblright{}, which CDIAL lists under Sanskrit \emph{bibheti}, \textquotedblleft{}fears\textquotedblright{} (CDIAL 9241). CDIAL does not list the noun itself.
\end{cousinweb}

\begin{grammarlens}[title={What you've built}]
One more word for how it feels.
\end{grammarlens}

\subsection*{Your turn}
\begin{itemize}
\raggedright
  \item \emph{Say it:} \emph{bīk}
  \item \emph{Say it:} \emph{bīk}, once more
  \item \emph{Say it:} say \emph{bīk}
  \item \emph{Recall:} say the Gujarati for happiness, then the Gujarati for anger, then say \emph{bīk} again
\end{itemize}

\begin{tcolorbox}[breakable,colback=teal!4,colframe=teal!35!black,title={Before you move on}]
What does \gu{બીક} mean? (\textquotedblleft{}Fear\textquotedblright{}.) Say it once more.
\end{tcolorbox}
\section[\texorpdfstring{\gu{ચિંતા} (cintā)}{ચિંતા (cintā)}]{\gu{ચિંતા} (cintā) — worry}

\label{lesson:GU-C50-worry}



\begin{quote}
Before the new one: say the Gujarati for anger, then the Gujarati for fear.
\end{quote}

\subsection*{You'll want to know: \gu{ચિંતા}}
\textbf{\gu{ચિંતા}} (\emph{cintā}) — \textquotedblleft{}worry\textquotedblright{}.

\begin{cousinweb}
\textbf{\gu{ચિંતા}} (\emph{cintā}) is Sanskrit \emph{cintā}, \textquotedblleft{}thought, care\textquotedblright{}, taken over whole as a learned loan, a tatsama (CDIAL 4816).
\end{cousinweb}

\begin{grammarlens}[title={What you've built}]
One more word for how it feels.
\end{grammarlens}

\subsection*{Your turn}
\begin{itemize}
\raggedright
  \item \emph{Say it:} \emph{cintā}
  \item \emph{Say it:} \emph{cintā}, once more
  \item \emph{Say it:} say \emph{cintā}
  \item \emph{Recall:} say the Gujarati for anger, then the Gujarati for fear, then say \emph{cintā} again
\end{itemize}

\begin{tcolorbox}[breakable,colback=teal!4,colframe=teal!35!black,title={Before you move on}]
What does \gu{ચિંતા} mean? (\textquotedblleft{}Worry\textquotedblright{}.) Say it once more.
\end{tcolorbox}
\section[\texorpdfstring{\gu{થાક} (thāk)}{થાક (thāk)}]{\gu{થાક} (thāk) — tiredness}

\label{lesson:GU-C50-tiredness}



\begin{quote}
Before the new one: say the Gujarati for fear, then the Gujarati for worry.
\end{quote}

\subsection*{You'll want to know: \gu{થાક}}
\textbf{\gu{થાક}} (\emph{thāk}) — \textquotedblleft{}tiredness\textquotedblright{}. What a long day leaves.

\begin{cousinweb}
\textbf{\gu{થાક}} (\emph{thāk}) belongs with \emph{thākvũ}, \textquotedblleft{}to become weary\textquotedblright{}, which CDIAL lists under a reconstructed \emph{*sthakk}, \textquotedblleft{}stop, halt\textquotedblright{} (CDIAL 13737). CDIAL thinks a link with \emph{sthā}, \textquotedblleft{}to stand\textquotedblright{}, probable but not clear.
\end{cousinweb}

\begin{grammarlens}[title={What you've built}]
One more word for how it feels.
\end{grammarlens}

\subsection*{Your turn}
\begin{itemize}
\raggedright
  \item \emph{Say it:} \emph{thāk}
  \item \emph{Say it:} \emph{thāk}, once more
  \item \emph{Say it:} say \emph{thāk}
  \item \emph{Recall:} say the Gujarati for fear, then the Gujarati for worry, then say \emph{thāk} again
\end{itemize}

\begin{tcolorbox}[breakable,colback=teal!4,colframe=teal!35!black,title={Before you move on}]
What does \gu{થાક} mean? (\textquotedblleft{}Tiredness\textquotedblright{}.) Say it once more.
\end{tcolorbox}
